\subsection{Paths and loops in a quiver}

DEFINITION 2. --- Let C be a quiver. A path in C is a sequence \(c = (a_{0}, f_{1}, a_{1}, \ldots, a_{n-1}, f_{n}, a_{n})\) , where \(n\) is an integer \(\geqslant 0\) , \(a_{0}, \ldots, a_{n}\) are vertices of C and where, for \(1 \leqslant i \leqslant n\) , \(f_{i}\) is an arrow of C connecting \(a_{i-1}\) to \(a_{i}\) . We say that the path c has length \(n\) .

Let \(c = (a_{0}, f_{1}, a_{1}, \ldots, a_{n-1}, f_{n}, a_{n})\) be a path of length n in C. The vertex \(a_{0}\) is called the origin, or source, of c; the vertex \(a_{n}\) is called the term, or the target, of c; one also says that c joins the vertex \(a_{0}\) to the vertex \(a_{n}\) . The vertex \(a_{i}\) , for \(0 \leqslant i \leqslant n\) , is called the vertex of index i; the arrow \(f_{i}\) , for \(1 \leqslant i \leqslant n\) , is called the i-th arrow, or the arrow of index i. A path of length \(\geqslant 1\) is determined by the sequence of its arrows.

A path whose origin is equal to its endpoint is called a loop. A path of length 0 is a loop, said to be constant.

We say that paths

\[
c = (a _ {0}, f _ {1}, a _ {1}, \ldots , a _ {n - 1}, f _ {n}, a _ {n}) \text{and} c ^ {\prime} = (a _ {0} ^ {\prime}, f _ {1} ^ {\prime}, a _ {1} ^ {\prime}, \ldots , a _ {m - 1} ^ {\prime}, f _ {m} ^ {\prime}, a _ {m} ^ {\prime})
\]

in C are juxtaposable if the endpoint \(a_{n}\) of c is the origin \(a_{0}^{\prime}\) of \(c^{\prime}\) . In this case, the sequence \((a_{0}, f_{1}, a_{1}, \ldots, a_{n-1}, f_{n}, a_{n}, f_{1}^{\prime}, a_{1}^{\prime}, \ldots, f_{m}^{\prime}, a_{m}^{\prime})\) is a path in C, denoted \(c * c^{\prime}\) and called the juxtaposed path of c and \(c^{\prime}\) . It joins the origin of c to the endpoint of \(c^{\prime}\) ; its length is the sum of the lengths of the paths c and \(c^{\prime}\) .

